\documentclass[12pt,a4paper]{article}
\usepackage[utf8]{inputenc}
\usepackage{geometry}
\geometry{a4paper, margin=1in}
\usepackage{titlesec}
\usepackage{hyperref}
\usepackage{graphicx}
\usepackage{listings}
\usepackage{color}

\definecolor{codegreen}{rgb}{0,0.6,0}
\definecolor{codegray}{rgb}{0.5,0.5,0.5}
\definecolor{codepurple}{rgb}{0.58,0,0.82}
\definecolor{backcolour}{rgb}{0.95,0.95,0.92}

\lstdefinestyle{mystyle}{
    backgroundcolor=\color{backcolour},   
    commentstyle=\color{codegreen},
    keywordstyle=\color{magenta},
    numberstyle=\tiny\color{codegray},
    stringstyle=\color{codepurple},
    basicstyle=\ttfamily\footnotesize,
    breakatwhitespace=false,         
    breaklines=true,                 
    captionpos=b,                    
    keepspaces=true,                 
    numbers=left,                    
    numbersep=5pt,                  
    showspaces=false,                
    showstringspaces=false,
    showtabs=false,                  
    tabsize=2
}
\lstset{style=mystyle}

\title{
    \vspace{-2cm}
    \textbf{EGC 301P: Operating Systems Lab} \\
    \large Mini-Project Report \\
    \vspace{0.5cm}
    \textbf{Competitive Programming Platform with 1v1 Duels}
}
\author{Vedansh Patel (BT2024162) \\ \href{https://github.com/username/repo}{github.com/username/repo}}
\date{}
\begin{document}

\maketitle

\section{Problem Statement}
The objective of this project was to design and implement a \textbf{Competitive Programming (CP) Platform} that allows users to solve programming problems, submit their solutions, and compete with other users in real time 1v1 duels, It also has an ELO system like chess where users can lose or gain rating and climb the leaderboard. The system evaluates the submitted code against predefined test cases within a safe environment (using various OS library functions to stop malicious users from gaining unautorized access to system resources).

This project was built in C, adhering to a client-server architecture, using Operating Systems concepts such as concurrency, inter process communication, file locking, and socket programming.

\section{Implementation of Mandatory OS Concepts}

\subsection{Role-Based Authorization}
The system has 2 distinct user roles to ensure proper role based access control. Roles are encoded as privileges within the 'UserRecord' structure:
\begin{itemize}
    \item \texttt{REGULAR\_USER (0)}: Can log in, view problems, submit solutions, participate in duels, view the leaderboard.
    \item \texttt{ADMIN\_USER (1)}: Has elevated privileges to manage the platform, such as managing database of problems.
\end{itemize}
During initialization ('serverInit()'), an admin user account is explicitly generated. Client requests are validated against their authenticated session's 'privilege' attribute, restricting sensitive operations like \texttt{ADD\_PROBLEM} or \texttt{DELETE\_PROBLEM} to only authorized roles.

\begin{lstlisting}[language=C, caption=UserRecord Structure for Authorization]
typedef struct{
    char username[100];
    char password[100];
    int privilege; // 0 = regular user, 1 = admin
}UserRecord;
\end{lstlisting}
\subsection{File Locking}
To maintain the integrity of the generic database (`users.dat`, `problems.dat`, `leaderboard.dat`), safe file access is ensured using advisory record locking via \texttt{fcntl()}.
\begin{itemize}
    \item \textbf{Read Locks (\texttt{F\_RDLCK})}: Applied when retrieving records (e.g., `getRecordCount`, `getRecord`) to allow concurrent reads but prevent writes.
    \item \textbf{Write Locks (\texttt{F\_WRLCK})}: Applied when inserting, updating, or deleting records (e.g., `insertRecord`, `updateRecord`) to ensure exclusive access.
    \item Operations utilize \texttt{F\_SETLKW} to block the calling thread until the lock is successfully acquired, guaranteeing safety without manual sleep loops.
\end{itemize}

\begin{lstlisting}[language=C, caption=File Locking Implementation]
struct flock getReadLock(off_t offset, off_t length){
    return getFileLock(F_RDLCK, offset, length);
}

// Applying the lock
struct flock lock = getReadLock(0, sizeof(int));
if(fcntl(fd, F_SETLKW, &lock) < 0){
    perror("failed to acquire read lock");
    close(fd);
    return -1;
}
\end{lstlisting}

\subsection{Concurrency Control}
Given the multi-threaded nature of the server handling multiple clients simultaneously, synchronization mechanisms are heavily used:
\begin{itemize}
    \item \textbf{Mutexes (\texttt{pthread\_mutex\_t})}: Used to protect shared structures. For instance, \texttt{activeUsersLock} safeguards operations on the \texttt{activeUserList} array, and \texttt{challengeLock} prevents data corruption when users issue or retract duel requests.
    \item \textbf{Semaphores (\texttt{sem\_t})}: A semaphore named \texttt{job\_queue} initialized to \texttt{MAX\_JOBS} (set to 2) is used in the `submitProblem` function. This limits the number of concurrent code compilation and execution jobs, preventing the system from being overwhelmed by a burst of simultaneous solution submissions (or an attack from malicious people).
\end{itemize}

\begin{lstlisting}[language=C, caption=Concurrency Control using Semaphores]
sem_t job_queue;
sem_init(&job_queue, 0, MAX_JOBS);

// Inside submitProblem
sem_wait(&job_queue);
// ... job compilation and execution ...
sem_post(&job_queue);
\end{lstlisting}

\subsection{Data Consistency}
Data consistency is actively maintained across the application to prevent anomalies like race conditions, dirty reads, and lost updates:
\begin{itemize}
    \item \textbf{Race Conditions}: Mutexes protect state changes (like adding/removing users from arrays). Condition variables (\texttt{pthread\_cond\_t duelCond}) are paired with user-specific mutexes to reliably wake up threads when their 1v1 duel states change (To prevent corrupted reads).
    \item \textbf{Dirty Reads/Lost Updates}: By locking database file chunks before performing any Read-Modify-Write cycle (such as updating Elos using `updateRecord`), the system ensures that user data and scores are not overridden by concurrent socket interactions.
\end{itemize}

\begin{lstlisting}[language=C, caption=Preventing Race Conditions with Mutexes]
int removeActiveUser(char* username){
    pthread_mutex_lock(&activeUsersLock);
    for(int i = 0; i < activeUsers; i++){
        if(activeUserList[i] != NULL && strcmp(activeUserList[i]->username, username) == 0){
            activeUserList[i] = activeUserList[activeUsers - 1];
            activeUserList[activeUsers - 1] = NULL;
            activeUsers--;
            pthread_mutex_unlock(&activeUsersLock);
            return 0;
        }
    }
    pthread_mutex_unlock(&activeUsersLock);
    return -1;
}
\end{lstlisting}

\subsection{Socket Programming}
The platform utilizes a robust TCP Client-Server Model:
\begin{itemize}
    \item The Server opens a passive listening socket and accepts incoming connections, spawning a new \texttt{pthread} (`handleClient()`) for each connected client.
    \item Communication relies heavily on a structured JSON protocol (using the `cJSON` library) sent over standard `read()` and `write()` calls.
    \item The `select()` system call is leveraged in the `startDuel()` routine to actively monitor input from both duel participants simultaneously, establishing a responsive real-time event loop without busy-waiting, and maintaining fairness to both participants.
\end{itemize}

\begin{lstlisting}[language=C, caption=Monitoring Multiple Sockets with Select]
fd_set readfds;
FD_ZERO(&readfds);
FD_SET(arena->socket1, &readfds);
FD_SET(arena->socket2, &readfds);

int activity = select(maxFd + 1, &readfds, NULL, NULL, NULL);

if(FD_ISSET(arena->socket1, &readfds)){
    int res = handleDuelSocketMessage(..., arena->socket1, ...);
}
\end{lstlisting}

\subsection{Inter-Process Communication (IPC)}
When evaluating user submissions, the server safely forks a child process to compile and run the arbitrary code. Pipes are heavily utilized for IPC:
\begin{itemize}
    \item Anonymous Pipes (\texttt{pipe()}) are instantiated before \texttt{fork()} to capture compilation logs and execution runtime outputs.
    \item The child process uses \texttt{dup2()} to redirect \texttt{STDERR\_FILENO} and \texttt{STDOUT\_FILENO} to the pipe's write end.
    \item The parent server thread uses a custom `readPipe()` function to seamlessly read the diagnostic outputs and provide meaningful error messages (like Compilation Errors or Sandboxing Violations) back to the client.
\end{itemize}

\begin{lstlisting}[language=C, caption=Using Pipes for IPC between Server and Compilation Sandbox]
int pipefd[2];
pipe(pipefd);

int pid = fork();
if(pid == 0){
    close(pipefd[0]);
    dup2(pipefd[1], STDERR_FILENO); // Redirect stderr to pipe
    close(pipefd[1]);
    
    execlp("gcc", "gcc", "-o", tempExecutablePath, tempSolutionPath, NULL);
    exit(-1);
}
close(pipefd[1]);
int bytesRead = readPipe(pipefd[0], buffer, sizeof(buffer), "Compilation error: ");
\end{lstlisting}


\section{Challenges Faced and Solutions}

\begin{itemize}
    \item \textbf{Challenge:} Arbitrary Code Execution and Security. \newline
    \textbf{Solution:} Running user submitted C code natively on the server is highly dangerous. To mitigate this, `seccomp-bpf` a Linux kernel feature that uses BPF bytecode to filter system calls, these filters were implemented to drop harmful system calls (like execve (to execute rm -rf *), or sleep(1000) to sleep forever and take up a slot), and `setrlimit` was configured to restrict memory (\texttt{RLIMIT\_AS}), CPU time (\texttt{RLIMIT\_CPU}), thus creating a safe environment for the code to execute without being harmful to the main server.
    
    \item \textbf{Challenge:} Handling Concurrent Duels without blocking main threads.\newline
    \textbf{Solution:} Instead of tying up client-handling threads in busy waits, the `select()` system call was utilized to monitor both participant's sockets simultaneously.
    
    \item \textbf{Challenge:} Managing Database Integrity across multiple reads/writes.\newline
    \textbf{Solution:} Since a relational database wasn't used, a custom generic file-based database utility (`genericDatabase.c`) was implemented that correctly applies `fcntl` advisory locks at the OS level to ensure safe interactions.
    
    \item \textbf{Challenge:} Code duplication and complexity when managing separate database stores (Leaderboard, User Data, Problems). \newline
    \textbf{Solution:} Refactored the database logic to create a highly modular generic database store. By utilizing function pointers to dynamically handle different structures, the need for separate implementations was eliminated. This drastically reduced the number of lines of code and significantly improved overall modularity.

    \item \textbf{Challenge:} Transferring large files (10 mB+) over the socket easily overflows fixed 4kB memory buffers, causing incomplete transmissions or data loss. \newline
    \textbf{Solution:} Designed a chunked data transmission protocol. First, the sender transmits the total file size as a header so the receiver knows exactly how much data to expect. The file is then iteratively broken down and sent in safe, manageable chunks. On the receiving end, a loop utilizing the 'read()' system call accumulates these chunks until the entire file is correctly reconstructed, completely avoiding buffer overflow.
    
\end{itemize}

\end{document}
